\documentclass[12pt]{article}
\usepackage[utf8]{inputenc}
\usepackage[T1]{fontenc}
\usepackage[margin=1in]{geometry}
\usepackage{setspace}
\usepackage{xcolor}
\usepackage[colorlinks=true, urlcolor=blue!60!black, linkcolor=blue!60!black]{hyperref}
\setlength{\parindent}{2em}
\setlength{\parskip}{1.25em}
\renewcommand{\baselinestretch}{1.0}
\begin{document}
\begin{center}
{\LARGE \textit{Korematsu v. United States}, 323 U.S. 214 (1944)\par}
\vspace{0.8em}
{\large Case Brief\par}
\end{center}
\vspace{1em}

\section{Memory Jogger}
Deferring to the combined wartime authority of the President and Congress, the Supreme Court upheld the criminal conviction of an American citizen of Japanese ancestry for defying a World War II West Coast exclusion order---a decision later repudiated as ``gravely wrong the day it was decided'' (\href{https://www.courtlistener.com/opinion/104040/korematsu-v-united-states/}{CourtListener}; \href{https://constitutioncenter.org/blog/on-this-day-the-supreme-court-issues-the-korematsu-decision}{National Constitution Center}).

\section{Detailed Case Facts}
Following the Japanese attack on Pearl Harbor on December 7, 1941, and the United States' declaration of war on December 8, 1941, President Franklin D. Roosevelt signed Executive Order 9066 on February 19, 1942 (\href{https://www.law.cornell.edu/supremecourt/text/323/214}{Cornell LII}; \href{https://constitutioncenter.org/blog/on-this-day-the-supreme-court-issues-the-korematsu-decision}{National Constitution Center}). Citing protection against espionage and sabotage, Executive Order 9066 authorized the Secretary of War and designated military commanders to prescribe military areas from which ``any or all persons may be excluded'' (\href{https://www.law.cornell.edu/supremecourt/text/323/214}{Cornell LII}). On February 20, 1942, Lieutenant General John L. DeWitt was designated Commanding General of the Western Defense Command, and on March 2, 1942, he issued Public Proclamation No.\ 1 designating the Pacific Coast states---including California---as Military Areas Nos.\ 1 and 2 (\href{https://www.law.cornell.edu/supremecourt/text/323/214}{Cornell LII}). On March 21, 1942, Congress enacted a statute (56 Stat.\ 173, 18 U.S.C.\ \S~97a) making it a federal misdemeanor, punishable by up to a \$5,000 fine and one year's imprisonment, knowingly to enter, remain in, leave, or violate restrictions in any prescribed military area (\href{https://www.law.cornell.edu/supremecourt/text/323/214}{Cornell LII}).

On March 27, 1942, General DeWitt issued Public Proclamation No.\ 4, prohibiting persons of Japanese ancestry from leaving Military Area No.\ 1 starting March 29, 1942, until directed by a future order (\href{https://www.law.cornell.edu/supremecourt/text/323/214}{Cornell LII}). Then, on May 3, 1942, DeWitt issued Civilian Exclusion Order No.\ 34, directing that all persons of Japanese ancestry, both alien and non-alien, be excluded from a portion of Military Area No.\ 1 including Alameda County, California, after May 9, 1942, and instructing them to report to a Civil Control Station for transfer to an Assembly Center (\href{https://www.law.cornell.edu/supremecourt/text/323/214}{Cornell LII}).

Fred Toyosaburo Korematsu, a native-born American citizen of Japanese descent whose loyalty to the United States was never questioned, lived and worked as a welder in San Leandro, Alameda County, California (\href{https://www.law.cornell.edu/supremecourt/text/323/214}{Cornell LII}; \href{https://www.oyez.org/cases/1940-1955/323us214}{Oyez}). Wishing to remain in his home community rather than submit to forced removal and detention, Korematsu stayed in San Leandro past the May 9 deadline and was arrested on May 30, 1942 (\href{https://www.law.cornell.edu/supremecourt/text/323/214}{Cornell LII}).

\section{Procedural History}
On June 12, 1942, federal prosecutors filed an information in the U.S. District Court for the Northern District of California charging Korematsu with knowingly remaining in the prohibited area in violation of the Act of March 21, 1942, and Civilian Exclusion Order No.\ 34 (\href{https://www.law.cornell.edu/supremecourt/text/323/214}{Cornell LII}). After his demurrer was overruled, Korematsu was convicted following a bench trial; the district court suspended imposition of sentence and placed him on five years' probation, after which military authorities immediately took him into custody and placed him in an Assembly Center and later a Relocation Center (\href{https://www.law.cornell.edu/supremecourt/text/323/214}{Cornell LII}; \href{https://supreme.justia.com/cases/federal/us/319/432/}{Justia, 319 U.S. 432}).

When Korematsu appealed, the Ninth Circuit certified a jurisdictional question to the Supreme Court, which held in \textit{Korematsu v. United States}, 319 U.S. 432 (1943), that a probation order without a formal prison sentence or fine is a final, appealable judgment (\href{https://supreme.justia.com/cases/federal/us/319/432/}{Justia, 319 U.S. 432}). On the merits, the Ninth Circuit affirmed his conviction, 140 F.2d 289 (9th Cir. 1943) (\href{https://law.justia.com/cases/federal/appellate-courts/F2/140/289/1567033/}{Justia, 140 F.2d 289}). The Supreme Court granted certiorari (321 U.S. 760), heard argument on October 11--12, 1944, and affirmed on December 18, 1944 (rehearing denied Feb.\ 12, 1945, 324 U.S. 885) (\href{https://www.law.cornell.edu/supremecourt/text/323/214}{Cornell LII}; \href{https://tile.loc.gov/storage-services/service/ll/usrep/usrep323/usrep323214/usrep323214.pdf}{Library of Congress PDF}).

\textbf{Subsequent History:} In 1984, after researchers uncovered suppressed Justice Department and naval intelligence documents showing that the government had misled the Supreme Court about the lack of military necessity, Judge Marilyn Hall Patel vacated Korematsu's conviction on a writ of error \textit{coram nobis} in \textit{Korematsu v. United States}, 584 F. Supp. 1406 (N.D. Cal. 1984) (\href{https://texaslawreview.org/korematsu-hirabayashi-and-the-second-monster/}{Texas Law Review}). Congress apologized and enacted reparations in the Civil Liberties Act of 1988, Acting Solicitor General Neal Katyal issued an official Confession of Error in 2011, and in \textit{Trump v. Hawaii}, 585 U.S. 667 (2018), Chief Justice Roberts wrote for the Court that \textit{Korematsu} ``was gravely wrong the day it was decided, has been overruled in the court of history, and---to be clear---`has no place in law under the Constitution' '' (\href{https://constitutioncenter.org/blog/on-this-day-the-supreme-court-issues-the-korematsu-decision}{National Constitution Center}; \href{https://www.justice.gov/archives/opa/blog/confession-error-solicitor-generals-mistakes-during-japanese-american-internment-cases}{DOJ Confession of Error}).

\section{Judicial Votes}
The Supreme Court decided the case by a \textbf{6--3} vote (\href{https://www.oyez.org/cases/1940-1955/323us214}{Oyez}; \href{https://www.law.cornell.edu/supremecourt/text/323/214}{Cornell LII}):
\begin{itemize}
    \item \textbf{Majority (6):} Justice Hugo Black authored the opinion of the Court, joined by Chief Justice Harlan Fiske Stone and Justices Stanley Reed, Felix Frankfurter, William O. Douglas, and Wiley B. Rutledge (\href{https://www.oyez.org/cases/1940-1955/323us214}{Oyez}).
    \item \textbf{Concurrence (1):} Justice Felix Frankfurter joined the majority opinion and filed a separate concurring opinion (\href{https://www.law.cornell.edu/supremecourt/text/323/214}{Cornell LII}).
    \item \textbf{Dissents (3):} Justice Owen J. Roberts, Justice Frank Murphy, and Justice Robert H. Jackson each filed separate dissenting opinions (\href{https://www.law.cornell.edu/supremecourt/text/323/214}{Cornell LII}).
\end{itemize}

\section{Holding}
\textbf{Judgment:} Affirmed; Korematsu's conviction under the Act of March 21, 1942, for violating Civilian Exclusion Order No.\ 34 was upheld (\href{https://www.law.cornell.edu/supremecourt/text/323/214}{Cornell LII}).

\textbf{Rule Announced:} Legal restrictions that curtail the civil rights of a single racial group are ``immediately suspect'' and must be subjected to ``the most rigid scrutiny,'' justifiable only by ``pressing public necessity'' and never by racial antagonism; nevertheless, during wartime, when Congress and the President act jointly under their constitutional war powers and military authorities apprehend ``the gravest imminent danger to the public safety,'' temporary exclusion of citizens of Japanese ancestry from a threatened military area is constitutionally permissible (\href{https://www.law.cornell.edu/supremecourt/text/323/214}{Cornell LII}; \href{https://constitutioncenter.org/blog/on-this-day-the-supreme-court-issues-the-korematsu-decision}{National Constitution Center}).

\textbf{Holding vs.\ Dicta:} The narrow holding was confined solely to the validity of Civilian Exclusion Order No.\ 34 as of May 1942, when Korematsu violated it, while explicitly refusing to decide the constitutionality of subsequent indeterminate detention in assembly or relocation centers---an issue addressed on statutory grounds the same day in \textit{Ex parte Endo}, 323 U.S. 283 (1944) (\href{https://www.law.cornell.edu/supremecourt/text/323/214}{Cornell LII}). Irony surrounds the opinion's most enduring passage: Justice Black's statement that racial classifications are ``immediately suspect'' and require ``the most rigid scrutiny'' operated as a doctrinal pronouncement that became the foundation of modern strict scrutiny even though the Court failed to apply meaningful scrutiny to the government's actual record (\href{https://constitutioncenter.org/blog/on-this-day-the-supreme-court-issues-the-korematsu-decision}{National Constitution Center}).

\section{Analysis of Opinions}

\subsection{Majority Opinion (Justice Black)}
Justice Black opened with the landmark principle that all legal restrictions curtailing the civil rights of a single racial group are ``immediately suspect'' and subject to ``the most rigid scrutiny'' (\href{https://www.law.cornell.edu/supremecourt/text/323/214}{Cornell LII}). Yet he immediately pivoted to wartime separation-of-powers deference, relying heavily on \textit{Hirabayashi v. United States}, 320 U.S. 81 (1943), which had upheld a West Coast curfew order under the same Executive Order 9066 and Act of March 21, 1942 (\href{https://www.law.cornell.edu/supremecourt/text/323/214}{Cornell LII}). Black reasoned that the Court could not ``reject as unfounded the judgment of the military authorities and of Congress'' that disloyal individuals existed within the Japanese-American population and could not be quickly separated from the loyal majority (\href{https://www.law.cornell.edu/supremecourt/text/323/214}{Cornell LII}).

Black rejected Korematsu's argument that Public Proclamation No.\ 4 (forbidding him to leave) and Exclusion Order No.\ 34 (forbidding him to remain) created a contradictory trap, reasoning that the May 3 exclusion order superseded the March 27 freeze (\href{https://www.law.cornell.edu/supremecourt/text/323/214}{Cornell LII}). Crucially, he severed the exclusion order from the detention program, insisting that because Korematsu had not been convicted of failing to report to or remain in an assembly or relocation center, the Court was passing only on exclusion (\href{https://www.law.cornell.edu/supremecourt/text/323/214}{Cornell LII}). Denying that the camps were ``concentration camps'' or that exclusion stemmed from racial hostility, Black concluded that the war-making branches acted out of military urgency and that the Court could not second-guess them ``by availing ourselves of the calm perspective of hindsight'' (\href{https://www.law.cornell.edu/supremecourt/text/323/214}{Cornell LII}).

\subsection{Concurring Opinion (Justice Frankfurter)}
Justice Frankfurter joined Black's opinion and emphasized the constitutional structure of the war power. Quoting Chief Justice Hughes, he stressed that the war power is ``the power to wage war successfully'' and is ``as much part of the Constitution as provisions looking to a nation at peace'' (\href{https://www.law.cornell.edu/supremecourt/text/323/214}{Cornell LII}). Directly answering Justice Jackson's dissent, Frankfurter rejected the notion that an order could be a valid military precaution yet constitutionally lawless, asserting that if a military order is within the war power, Congress may enforce it in civil courts, while adding that holding a measure constitutional ``does not carry with it approval of that which Congress and the Executive did. That is their business, not ours'' (\href{https://www.law.cornell.edu/supremecourt/text/323/214}{Cornell LII}).

\subsection{Dissenting Opinion (Justice Roberts)}
Justice Roberts dismantled the majority's attempt to sever exclusion from detention. Tracing the exact chronology of Public Proclamation No.\ 4 and Exclusion Order No.\ 34, he demonstrated that Korematsu faced ``two diametrically contradictory orders'': he committed a crime if he left Military Area No.\ 1 and a crime if he stayed anywhere outside an Assembly Center---a ``cleverly devised trap'' whose sole purpose was ``imprisonment in a concentration camp, based on his ancestry, and solely because of his ancestry, without evidence or inquiry concerning his loyalty'' (\href{https://www.law.cornell.edu/supremecourt/text/323/214}{Cornell LII}). Punishing a citizen under such conflicting commands violated Fifth Amendment due process (\href{https://www.law.cornell.edu/supremecourt/text/323/214}{Cornell LII}).

\subsection{Dissenting Opinion (Justice Murphy)}
Justice Murphy condemned the exclusion as going ``over `the very brink of constitutional power' '' and falling ``into the ugly abyss of racism'' (\href{https://www.law.cornell.edu/supremecourt/text/323/214}{Cornell LII}). While acknowledging wide military discretion during war, he insisted that in the absence of martial law, courts must verify whether a deprivation of rights is reasonably related to an ``immediate, imminent, and impending'' public danger (\href{https://www.law.cornell.edu/supremecourt/text/323/214}{Cornell LII}). Dissecting General DeWitt's Final Report---which branded Japanese Americans ``an enemy race''---Murphy showed that the evacuation rested on discredited racial stereotypes rather than evidence, noting that not a single person of Japanese ancestry had been accused or convicted of espionage or sabotage after Pearl Harbor, that months elapsed before removal, and that individualized loyalty hearings (like those used for German and Italian aliens and by Britain for 74,000 enemy aliens) were entirely feasible (\href{https://www.law.cornell.edu/supremecourt/text/323/214}{Cornell LII}).

\subsection{Dissenting Opinion (Justice Jackson)}
Justice Jackson offered a profound separation-of-powers critique. He first noted that Korematsu was convicted solely for being present in his home state because of his ancestry, violating the constitutional principle that ``guilt is personal and not inheritable'' (\href{https://www.law.cornell.edu/supremecourt/text/323/214}{Cornell LII}). Turning to the institutional role of courts, Jackson argued that civil judges lack the evidence and competence to evaluate military necessity and must not ``distort the Constitution to approve all that the military may deem expedient'' (\href{https://www.law.cornell.edu/supremecourt/text/323/214}{Cornell LII}). A temporary military overstep is an ``incident,'' he warned, but when the Supreme Court rationalizes it as constitutional, the resulting principle ``lies about like a loaded weapon ready for the hand of any authority that can bring forward a plausible claim of an urgent need'' (\href{https://www.law.cornell.edu/supremecourt/text/323/214}{Cornell LII}). Civil courts should therefore refuse to enforce unconstitutional military orders rather than become ``instruments of military policy'' (\href{https://www.law.cornell.edu/supremecourt/text/323/214}{Cornell LII}).

\section{Examples: Future Applications}
\textit{Note: Because \textit{Korematsu} was formally repudiated in \textit{Trump v. Hawaii}, 585 U.S. 667 (2018), the applications below analyze how modern separation-of-powers doctrine treats wartime and national-security executive actions in light of \textit{Korematsu}'s repudiation as an anticanonical negative precedent.}
\begin{enumerate}
    \item \textbf{Same-Side (Invalidated Group Detention of Citizens):} During a military conflict abroad, the President issues an executive order, later backed by a general congressional appropriation, requiring all U.S. citizens of a particular national origin living near coastal military bases to report to federal holding centers without individualized hearings. Under the repudiation of \textit{Korematsu} in \textit{Trump v. Hawaii} (and \textit{Ex parte Endo}), this is ``objectively unlawful and outside the scope of Presidential authority'' (\href{https://constitutioncenter.org/blog/on-this-day-the-supreme-court-issues-the-korematsu-decision}{National Constitution Center}).
    \item \textbf{Same-Side (Unchecked Executive Detention Without Due Process):} Federal officials detain a U.S. citizen domestically without charge based solely on an unsworn military commander's assertion of ``pressing public necessity,'' arguing courts must defer to executive security judgments. Following \textit{Hamdi v. Rumsfeld}, 542 U.S. 507 (2004), and \textit{Korematsu}'s rejection, courts refuse to rubber-stamp executive detention of citizens without meaningful notice and an opportunity to be heard (\href{https://supreme.justia.com/cases/federal/us/542/507/}{Justia, Hamdi v. Rumsfeld}).
    \item \textbf{Opposite-Side (Facially Neutral Foreign-Entry Restrictions):} Pursuant to explicit statutory authority under the Immigration and Nationality Act, the President suspends entry of foreign nationals from specified countries following a multi-agency security review of foreign vetting protocols. In \textit{Trump v. Hawaii}, 585 U.S. 667 (2018), the Supreme Court distinguished \textit{Korematsu} and upheld such a policy under deferential review because it involved foreign nationals seeking admission rather than the forcible relocation of U.S. citizens based explicitly on race (\href{https://constitutioncenter.org/blog/on-this-day-the-supreme-court-issues-the-korematsu-decision}{National Constitution Center}).
    \item \textbf{Opposite-Side (Temporary, Non-Racial Emergency Evacuation):} During an active coastal military bombardment or natural disaster, local and military authorities temporarily evacuate all residents of every race from a two-mile blast zone for 48 hours and allow them to travel freely elsewhere in the country. Even Justice Roberts's \textit{Korematsu} dissent recognized that temporary, non-discriminatory emergency exclusions (akin to drawing fire lines) are constitutionally permissible (\href{https://www.law.cornell.edu/supremecourt/text/323/214}{Cornell LII}).
    \item \textbf{Fence-Sitter (Targeted National-Security Surveillance or Screening):} During a declared international crisis, Congress authorizes an executive agency to apply heightened airport screening and financial-transaction audits to U.S. citizens who have recently traveled to a hostile nation, disproportionately affecting one ethnic community. Challengers argue this revives \textit{Korematsu}'s group-suspicion logic without individualized probable cause, whereas the government argues that travel-based, facially neutral criteria authorized by Congress merit substantial separation-of-powers deference short of detention (\href{https://www.scotusblog.com/2018/06/symposium-trump-v-hawaii-korematsus-ghost-and-national-security-masquerades/}{SCOTUSblog Symposium}).
\end{enumerate}

\section{Critique}
\textbf{Originalist and Textualist Perspectives:} Originalist and textualist scholars universally place \textit{Korematsu} alongside \textit{Dred Scott} and \textit{Plessy} in the constitutional ``anticanon'' (\href{https://pacificlegal.org/originalism-vs-textualism-vs-living-constitutionalism/}{Pacific Legal Foundation}). From a textualist and structural separation-of-powers standpoint, the Constitution's Fifth Amendment Due Process Clause, Sixth Amendment jury-trial guarantee, and Article I Suspension Clause forbid the executive branch from detaining U.S. citizens outside the criminal process unless Congress validly suspends the writ of habeas corpus in cases of rebellion or invasion (\href{https://www.californialawreview.org/print/constructive-constitutional-history-and-habeas-corpus-today}{California Law Review}). Originalists also fault the Court for permitting a sweeping, standardless congressional delegation to military commanders and for abandoning constitutional text in favor of unverified wartime pragmatism (\href{https://pacificlegal.org/originalism-vs-textualism-vs-living-constitutionalism/}{Pacific Legal Foundation}).

\textbf{Progressive and Doctrinal Perspectives:} Progressive and civil-rights scholars emphasize both the racial animus underlying Executive Order 9066 and the institutional failure of judicial review when executive officers withhold adverse evidence---a failure underscored by the 1984 \textit{coram nobis} vacatur and the Justice Department's 2011 Confession of Error (\href{https://www.justice.gov/archives/opa/blog/confession-error-solicitor-generals-mistakes-during-japanese-american-internment-cases}{DOJ Confession of Error}; \href{https://texaslawreview.org/korematsu-hirabayashi-and-the-second-monster/}{Texas Law Review}). Scholars such as Harold Hongju Koh further argue that even while formally repudiating \textit{Korematsu} in \textit{Trump v. Hawaii}, the Court risked resurrecting its ``ghost'' by granting excessive deference to executive national-security assertions in the face of evidence of group animus (\href{https://www.scotusblog.com/2018/06/symposium-trump-v-hawaii-korematsus-ghost-and-national-security-masquerades/}{SCOTUSblog Symposium}).

\textbf{Underlying Competing Values:} \textit{Korematsu} crystallizes the enduring separation-of-powers dilemma between (1) \textit{executive and legislative agility} in defending the nation during existential wartime crises, where political branches possess superior intelligence and operational responsibility, and (2) \textit{judicial independence and individual constitutional rights}, which exist precisely to check majoritarian panic and prevent the ``loaded weapon'' of unchecked emergency power from eroding the rule of law (\href{https://www.californialawreview.org/print/constructive-constitutional-history-and-habeas-corpus-today}{California Law Review}; \href{https://www.law.cornell.edu/supremecourt/text/323/214}{Cornell LII}).

\section{Key Quotations}
\begin{quote}
``It should be noted, to begin with, that all legal restrictions which curtail the civil rights of a single racial group are immediately suspect. That is not to say that all such restrictions are unconstitutional. It is to say that courts must subject them to the most rigid scrutiny. Pressing public necessity may sometimes justify the existence of such restrictions; racial antagonism never can.''
\end{quote}
\noindent --- Justice Hugo Black, Opinion of the Court (\href{https://constitutioncenter.org/blog/on-this-day-the-supreme-court-issues-the-korematsu-decision}{National Constitution Center}; \href{https://www.courtlistener.com/opinion/104040/korematsu-v-united-states/}{CourtListener}).

\begin{quote}
``Compulsory exclusion of large groups of citizens from their homes, except under circumstances of direst emergency and peril, is inconsistent with our basic governmental institutions. But when under conditions of modern warfare our shores are threatened by hostile forces, the power to protect must be commensurate with the threatened danger.''
\end{quote}
\noindent --- Justice Hugo Black, Opinion of the Court (\href{https://www.law.cornell.edu/supremecourt/text/323/214}{Cornell LII}).

\begin{quote}
``On the contrary, it is the case of convicting a citizen as a punishment for not submitting to imprisonment in a concentration camp, based on his ancestry, and solely because of his ancestry, without evidence or inquiry concerning his loyalty and good disposition towards the United States.''
\end{quote}
\noindent --- Justice Owen J. Roberts, Dissenting Opinion (\href{https://www.law.cornell.edu/supremecourt/text/323/214}{Cornell LII}).

\begin{quote}
``This exclusion of `all persons of Japanese ancestry, both alien and non-alien,' from the Pacific Coast area on a plea of military necessity in the absence of martial law ought not to be approved. Such exclusion goes over `the very brink of constitutional power' and falls into the ugly abyss of racism.''
\end{quote}
\noindent --- Justice Frank Murphy, Dissenting Opinion (\href{https://www.law.cornell.edu/supremecourt/text/323/214}{Cornell LII}).

\begin{quote}
``But once a judicial opinion rationalizes such an order to show that it conforms to the Constitution, or rather rationalizes the Constitution to show that the Constitution sanctions such an order, the Court for all time has validated the principle of racial discrimination in criminal procedure and of transplanting American citizens. The principle then lies about like a loaded weapon ready for the hand of any authority that can bring forward a plausible claim of an urgent need.''
\end{quote}
\noindent --- Justice Robert H. Jackson, Dissenting Opinion (\href{https://www.law.cornell.edu/supremecourt/text/323/214}{Cornell LII}).


\end{document}
